\section{Cardinal arithmetic}\label{sec:cardinal_arithmetic}

\paragraph{Cardinal arithmetic}

\begin{definition}\label{def:cardinal_arithmetic}
  We will define arithmetic operations for them. Unlike in \cref{def:ordinal_arithmetic}, we will directly define the operations as \hyperref[def:cardinality]{cardinal numbers} of some sets rather than via some form of recursion.

  Fix two ordinals \( \kappa \) and \( \mu \).
  \begin{thmenum}
    \thmitem{def:cardinal_arithmetic/addition}\mcite[def. 6.4.29]{Hinman2005Logic} Based on \cref{thm:ordinal_addition_disjoin_union}, we define their \term{sum} as
    \begin{equation*}
      \kappa + \mu \coloneqq \card(\kappa \amalg \mu),
    \end{equation*}
    where \( \kappa \amalg \mu \) is their \hyperref[def:disjoint_union]{disjoint union}.

    \thmitem{def:cardinal_arithmetic/multiplication}\mcite[def. 6.4.29]{Hinman2005Logic} Based on \cref{thm:ordinal_multiplication_cartesian_product}, we define their \term{product} as
    \begin{equation*}
      \kappa \cdot \mu \coloneqq \card(\kappa \times \mu).
    \end{equation*}

    \thmitem{def:cardinal_arithmetic/exponentiation}\mcite[exerc. 6.4.46]{Hinman2005Logic} We define \term{exponentiation} as
    \begin{equation*}
      \kappa^\mu \coloneqq \card(\fun(\kappa, \mu)).
    \end{equation*}
  \end{thmenum}
\end{definition}

\begin{proposition}\label{thm:cardinal_addition_algebraic_properties}
  Cardinal number addition is \hyperref[def:binary_operation/associative]{associative}, \hyperref[def:binary_operation/associative]{commutative} and \hyperref[def:binary_operation/cancellative]{cancellative}.

  Compare this with \cref{thm:natural_number_addition_properties} and \cref{thm:ordinal_addition_algebraic_properties}.
\end{proposition}
\begin{proof}
  Associativity and left cancellation is inherited from the ordinals. Commutativity and right cancellation hold because we are considering arbitrary bijective functions rather than the more restrictive order isomorphisms. Indeed, \( \kappa \amalg \mu \) and \( \mu \amalg \kappa \) may have different order types as demonstrated in \cref{ex:ordinal_addition_commutativity}, however there is an obvious bijective function between them.
\end{proof}

\begin{proposition}\label{thm:cardinal_multiplication_algebraic_properties}
  Cardinal number multiplication is \hyperref[def:binary_operation/associative]{associative}, \hyperref[def:binary_operation/associative]{commutative} and \hyperref[def:binary_operation/cancellative]{cancellative}.

  Compare this with \cref{thm:natural_number_multiplication_properties} and \cref{thm:ordinal_addition_algebraic_properties}.
\end{proposition}
\begin{proof}
  The result follows from the same considerations as in \cref{thm:cardinal_addition_algebraic_properties}.
\end{proof}

\begin{proposition}\label{thm:double_and_square_of_cardinal}
  For every cardinal \( \kappa \) we have
  \begin{thmenum}
    \thmitem{thm:double_and_square_of_cardinal/double} \( \kappa + \kappa = 2\kappa \).
    \thmitem{thm:double_and_square_of_cardinal/square} \( \kappa \cdot \kappa = \kappa^2 \).
  \end{thmenum}
\end{proposition}
\begin{proof}
  \SubProofOf{thm:double_and_square_of_cardinal/double} Obviously \( \kappa \amalg \kappa = 2 \times \kappa \).
  \SubProofOf{thm:double_and_square_of_cardinal/square} The function
  \begin{equation*}
    \begin{aligned}
      &T: \fun(2, \kappa) \to \kappa \times \kappa \\
      &T(f) \coloneqq (f(0), f(1))
    \end{aligned}
  \end{equation*}
  is clearly injective. It is also surjective because for any ordered pair \( (\gamma, \delta) \in \kappa \times \kappa \) we can define the function
  \begin{equation*}
    \begin{aligned}
      &f: 2 \to \kappa \\
      &f(k) \coloneqq \begin{cases}
        \gamma, k = 0 \\
        \delta, k = 1
      \end{cases}
    \end{aligned}
  \end{equation*}

  Then \( T(f) = (\gamma, \delta) \).
\end{proof}

\begin{proposition}\label{thm:simplified_cardinal_arithmetic}
  Unlike \hyperref[def:ordinal_arithmetic]{ordinal arithmetic} with its intricacies like \cref{ex:ordinal_addition_commutativity}, \hyperref[def:cardinal_arithmetic]{cardinal arithmetic} has a simpler behavior:
  \begin{thmenum}
    \thmitem{thm:simplified_cardinal_arithmetic/finite} If \( \kappa \) and \( \mu \) are finite cardinals, then \( \kappa + \mu \) and \( \kappa \cdot \mu \) are the familiar operations on \hyperref[def:natural_numbers]{natural numbers}.

    \thmitem{thm:simplified_cardinal_arithmetic/infinite}\mcite[6.4.35]{Hinman2005Logic} If either \( \kappa \) or \( \mu \) is infinite, then
    \begin{equation*}
      \kappa + \mu = \max\set{ \kappa, \mu }.
    \end{equation*}

    If, additionally, both are nonzero, then
    \begin{equation*}
      \kappa \cdot \mu = \max\set{ \kappa, \mu }.
    \end{equation*}
  \end{thmenum}
\end{proposition}
\begin{proof}
  \SubProofOf{thm:simplified_cardinal_arithmetic/finite} The addition of ordinals defined in \cref{def:ordinal_addition} is an extension of addition of natural numbers, hence the two are equivalent for finite ordinals. The equivalence with cardinal addition defined in \cref{def:cardinal_arithmetic/addition} comes from \cref{thm:ordinal_addition_disjoin_union} and the fact that every finite ordinal is a cardinal as demonstrated in \cref{thm:members_of_omega_are_cardinals}.

  Analogously, equivalence of cardinal and ordinal multiplication follows from \cref{thm:ordinal_multiplication_cartesian_product}.

  \SubProofOf{thm:simplified_cardinal_arithmetic/infinite} Suppose that either \( \kappa \) or \( \mu \) is infinite and let \( \nu \coloneqq \max\set{ \kappa, \mu } \). The cases where either of them is zero are trivial, hence suppose that both are nonzero.

  We have
  \begin{equation*}
    \kappa \amalg \mu \subseteq \nu \amalg \nu = 2 \amalg \nu \subseteq \nu \times \nu,
  \end{equation*}
  hence \( \kappa + \mu \leq \nu \cdot \nu = \nu^2 \).

  Furthermore there exists an obvious injective function from \( \nu \) to \( \kappa \amalg \mu \) (which is different depending on whether \( \nu = \kappa \) or \( \nu = \mu \)). Therefore,
  \begin{equation*}
    \nu \leq \kappa + \mu \leq \nu^2.
  \end{equation*}

  For multiplication we have \( \kappa \times \mu \subseteq \nu \times \nu \), hence
  \begin{equation*}
    \nu \leq \kappa \cdot \mu \leq \nu^2.
  \end{equation*}

  The rest follows from \cref{thm:square_of_infinite_ordinal}.
\end{proof}

\begin{corollary}\label{thm:finite_cardinal_exponent}
  For any infinite cardinal \( \kappa \), we have \( \kappa^n = \kappa \) for every positive natural number \( n \).
\end{corollary}
\begin{proof}
  Induction on \( n \), using \cref{thm:simplified_cardinal_arithmetic/infinite} at each step.
\end{proof}

\begin{definition}\label{def:strong_limit_cardinal}
\end{definition}

\begin{corollary}\label{thm:aleph_zero_is_strong_limit}
  The first infinite cardinal \( \aleph_0 \) is a \hyperref[def:strong_limit_cardinal]{strong limit cardinal}.
\end{corollary}
\begin{comments}
  \item See also \cref{thm:aleph_zero_is_regular}
\end{comments}
\begin{proof}
  \Cref{thm:simplified_cardinal_arithmetic/finite} states that cardinal exponentiation extends natural number exponentiation. Hence, we can conclude that \( 2^n < \aleph_0 \) for any \( n < \aleph_0 \) since the former are finite and the latter is not.
\end{proof}

\begin{proposition}\label{thm:cardinality_of_disjoint_union}
  If \( \mscrK \) has cardinality \( \kappa \) an if \( \mu \) bounds the cardinality of \( A_k \) for all indices \( k \), then the \hyperref[def:disjoint_union]{disjoint union} \( \coprod_{k \in \mscrK} A_k \) has cardinality at most \( \kappa \cdot \mu \).

  If the cardinality of \( A_k \) is exactly \( \mu \) for every index \( k \), then \( \coprod_{k \in \mscrK} A_k \) has cardinality \( \kappa \cdot \mu \).
\end{proposition}
\begin{comments}
  \item This generalizes the familiar characterization of multiplication as repeated addition.
\end{comments}
\begin{proof}
  By the definition of cardinality, there exist bijective functions \( \iota_\mscrK: \mscrK \to \kappa \) and \( \iota_{A_k}: A_k \to \mu \) for every \( k \) in \( \mscrK \). In case where the cardinality of \( A_k \) is only bounded by \( \mu \), the function \( \iota_{A_k} \) is injective but not necessarily surjective.

  Consider the function
  \begin{equation*}
    \begin{aligned}
      &\Phi:\coprod_{k \in \mscrK} A_k \to \aleph_0 \times \kappa, \\
      &\Phi(k, a_k) \coloneqq (\iota_\mscrK(k), \iota_{A_k}(a_k)).
    \end{aligned}
  \end{equation*}

  It is injective. Indeed, if \( \Phi(k, a_k) = \Phi(m, a_m) \), then \( k = m \) since \( \iota_\mscrK \) is injective, and \( a_k = a_m \) since \( \iota_{A_k} \) is injective. This shows that \( \Phi \) is injective.

  Therefore, \( \coprod_{k \in \mscrK} A_k \) is bounded by \( \aleph_0 \times \kappa \), which is in turn equinumerous with
  \begin{equation*}
    \kappa \cdot \mu = \card(\aleph_0 \times \kappa).
  \end{equation*}

  In the case where \( A_k \) has cardinality exactly \( \mu \) for every \( k \), \( \Phi \) is also surjective. Indeed, fix any natural number \( n \) and any ordinal \( \alpha \) in \( \kappa \). Let \( k \) be the preimage \( \iota_\mscrK^{-1}(n) \) and let \( a_k \) be the preimage \( \iota_{A_k}^{-1}(\alpha) \). Then \( \Phi(k, a_k) = (n, \alpha) \).

  In this case \( \coprod_{k \in \mscrK} A_k \) is equinumerous with \( \kappa \cdot \mu \).
\end{proof}

\begin{lemma}\label{thm:power_set_via_subsets}
  Fix a set \( A \). The power set \( \pow(A) \) is equinumerous with the set of \hyperref[def:boolean_function]{Boolean-valued functions} \( \pow(A, \set{ T, F }) \).

  More precisely, then the operator
  \begin{equation*}
    \begin{aligned}
      &T: \fun(A, \set{ T, F }) \to \pow(A) \\
      &T(f) \coloneqq \set{ f(x) = T \given x \in A }.
    \end{aligned}
  \end{equation*}
  is bijective.
\end{lemma}
\begin{proof}
  Injectivity is clear. To see surjectivity, fix some subset \( B \subset A \) and define
  \begin{equation*}
    \begin{aligned}
      &f: A \to \set{ T, F } \\
      &f(x) \coloneqq \begin{cases}
        T, &x \in B \\
        F, \T{otherwise}.
      \end{cases}
    \end{aligned}
  \end{equation*}

  Clearly \( f \in \fun(A, \set{ T, F }) \) and \( T(f) = B \).
\end{proof}

\begin{proposition}\label{thm:cardinal_exponentiation_power_set}
  For every set \( A \) we have
  \begin{equation*}
    \card(\pow(A)) = 2^{\card(A)}.
  \end{equation*}
\end{proposition}
\begin{proof}
  We can prove the proposition similarly to \cref{thm:power_set_via_subsets}, but much more convoluted.

  Let \( \varphi: A \to \card(A) \) be \( \psi: \pow(A) \to \card(\pow(A)) \) bijective functions. Note that
  \begin{equation*}
    2^{\card(A)} = \card(\fun(\card(A), \set{ 0, 1 }))
  \end{equation*}
  by definition. Let \( \theta: \fun(\card(A), \set{ 0, 1 }) \to 2^{\card(A)} \) be a bijective function.

  Define the operator
  \begin{equation*}
    \begin{aligned}
      &T: 2^{\card(A)} \to \card(\pow(A)) \\
      &T(p) \coloneqq \psi\parens{ \set{ \varphi^{-1}(\gamma) \given \gamma \in \card(A) \T{and} \theta^{-1}(p)(\gamma) = 1 } }.
    \end{aligned}
  \end{equation*}

  This operator is bijective since for any \( \delta \in \card(\pow(A)) \) we can define
  \begin{equation*}
    \begin{aligned}
      &f: \card(A) \to \set{ 0, 1 } \\
      &f(\gamma) \coloneqq \begin{cases}
        1, &\varphi^{-1}(\gamma) \in \psi^{-1}(\delta) \\
        0, &\T{otherwise.}
      \end{cases}
    \end{aligned}
  \end{equation*}
  so that \( T(\theta(f)) = \delta \).

  Therefore, \( \card(A) \) and \( \card(\pow(A)) \) are equinumerous.
\end{proof}

\begin{corollary}\label{thm:cardinal_exponentiation_comparison}
  A set is \hyperref[def:aleph_zero_related_cardinalities]{finite} if and only if its \hyperref[def:power_set]{power set} is finite.
\end{corollary}
\begin{proof}
  \SufficiencySubProof If \( A \) has cardinality \( n < \aleph_0 \), \cref{thm:cardinal_exponentiation_power_set} implies that the cardinality of \( \pow(A) \) is \( 2^n \). \Cref{thm:simplified_cardinal_arithmetic/finite} implies that \( 2^n \) coincides with natural number exponentiation, and is hence a finite cardinal.

  \NecessitySubProof Conversely, if \( \pow(A) \) is finite, since \( \pow(A) \) obviously dominates \( A \), then \( A \) is also finite.
\end{proof}

\begin{corollary}\label{thm:cardinal_exponentiation_comparison}
  We have \( \card A < \card B \) if and only if \( 2^{\card A} < 2^{\card B} \).
\end{corollary}
\begin{proof}
  \SufficiencySubProof Suppose that \( \card A < \card B \).

  Let \( f: A \to B \) be an injective function. It is not surjective by assumption. Then \( f(A) \subsetneq B \).

  Clearly \( \pow(f(A)) \subsetneq \pow(B) \). \Cref{thm:cardinal_exponentiation_power_set} then implies that \( 2^{\card A} < 2^{\card B} \).

  \NecessitySubProof Suppose that \( 2^{\card A} < 2^{\card B} \).
  \begin{itemize}
    \item If \( \card A = \card B \), then \( 2^{\card A} = 2^{\card B} \), which is a contradiction.
    \item If \( \card A > \card B \), then we have shown that \( 2^{\card A} > 2^{\card B} \), which is again a contradiction.
    \item It remains for \( \card A \) to be strictly less than \( \card B \).
  \end{itemize}
\end{proof}

\begin{proposition}\label{thm:strong_limit_cardinal_is_weak_limit}
  If \( \mu \) is the successor cardinal of \( \kappa \), then \( \mu \leq 2^\kappa \).

  In particular, every \hyperref[def:strong_limit_cardinal]{strong limit cardinal} is a \hyperref[def:strong_limit_cardinal]{limit cardinal}.

  Furthermore, if \( \kappa \) is infinite, then \fullref{hyp:generalized_continuum_hypothesis} implies that \( \mu = 2^\kappa \).
\end{proposition}
\begin{proof}
  It is clear from \fullref{thm:cantors_power_set_theorem} and \cref{thm:cardinal_exponentiation_power_set} that \( \kappa < 2^\kappa \). By definition, \( \mu \) is the smallest cardinal such that \( \kappa < \mu \). Therefore, \( \mu \leq 2^\kappa \).
\end{proof}
